\begin{abstract}
Recent developments exemplified by BitNet~\cite{bitnet} are ushering in a transformative phase characterized by 1-bit Large Language Models (LLMs). In this study, we propose a 1-bit LLM variant named \textbf{\text{BitNet b1.58}}, in which each parameter (weight) is represented in a ternary format, specifically chosen from \{-1, 0, 1\}. Remarkably, this model achieves parity with its full-precision (FP16 or BF16) Transformer LLM counterparts of equal model size and with equivalent training data, both in perplexity and in downstream task metrics. At the same time, it delivers pronounced improvements in efficiency—offering lower latency, reduced memory footprint, higher throughput, and diminished energy demands. Importantly, the advent of the 1.58-bit LLM establishes a novel scaling principle and provides actionable guidelines for training future LLMs that maintain both robust performance and cost-effectiveness. Additionally, this advancement paves the way for a reimagined computational framework and catalyzes opportunities for hardware specifically tailored for 1-bit LLM operations.
\end{abstract}

\section{The Era of 1-bit LLMs}

The recent trajectory of AI research has been marked by the rapid expansion in both scale and capability of Large Language Models (LLMs). These models exhibit exceptional proficiency across various natural language processing benchmarks. Nonetheless, their ever-growing parameter counts complicate practical deployment and intensify concerns related to resource usage, including heightened energy expenditure and associated economic or ecological impacts.

A promising direction to alleviate these issues involves leveraging post-training quantization to develop low-bit models optimized for inference~\cite{smoothquant, gptq, quip, quip_sharp}. By reducing the bit-width for both weights and activations, this approach can dramatically lower the computational and memory demands of LLMs. The field has progressed from utilizing 16-bit to even more compact representations, such as 4-bit solutions~\cite{gptq, awq}. However, despite their broad adoption in commercial settings, post-training quantization schemes are not without shortcomings, frequently falling short of optimality.

Recent advancements in 1-bit model designs, exemplified by BitNet~\cite{bitnet}, have opened up promising avenues for minimizing the computational expense of LLMs without significant loss in efficacy. Standard LLMs generally rely on 16-bit floating-point representations (FP16 or BF16), with matrix multiplications dominating the computational load. As a result, the primary source of computational overhead lies in the floating-point addition and multiplication processes. In sharp contrast, BitNet’s architecture confines matrix multiplications to integer addition, thereby reducing the energy requirements of LLMs by several orders of magnitude. Since power is a principal limiting factor for computational performance across numerous hardware platforms, these energy savings also facilitate increased processing speeds.

Beyond computation, the transfer of model parameters from DRAM to on-chip accelerators’ memory (e.g., SRAM) during inference constitutes another significant expense. Although expanding on-chip SRAM can raise throughput, this solution is often cost-prohibitive when compared to DRAM. 1-bit LLMs offer notable advantages over full-precision counterparts by drastically reducing memory usage—both in terms of total capacity and bandwidth needs. This reduction can greatly diminish the expense and duration of weight transfers from DRAM, thus enhancing inference efficiency and speed.

Here, we present a substantial innovation in the form of a 1-bit LLM variant named \textbf{\text{BitNet b1.58}}, where each parameter is restricted to the ternary set \{-1, 0, 1\}. By augmenting the initial 1-bit BitNet with a zero state, we arrive at an effective 1.58 bits per parameter. \text{BitNet b1.58} preserves all advantages of the original BitNet, such as the altered computational approach that eliminates nearly all multiplications during matrix multiplications and allows for substantial optimization. It also matches the original BitNet's energy efficiency and vastly surpasses FP16 LLMs with respect to memory consumption, throughput, and latency. Additionally, \text{BitNet b1.58} brings two novel strengths: first, its modeling power is enhanced via explicit feature filtering enabled by the introduction of zeros into the weight space, which can lead to notable performance gains for 1-bit LLMs. Second, empirical results indicate that \text{BitNet b1.58} can achieve parity with full-precision (FP16) models—matching both perplexity and end-task outcomes—at the 3B parameter scale under equivalent configurations (including model size and training data).

\section{BitNet b1.58}

\text{BitNet b1.58} builds upon the BitNet framework, which modifies the standard Transformer by replacing the \emph{nn.Linear} layers with \emph{BitLinear}. The model is trained from the ground up with weights quantized to 1.58 bits and activations to 8 bits. This iteration introduces several adjustments relative to the original BitNet, outlined below.

\paragraph{Quantization Function.} To ensure weight values are restricted to -1, 0, or +1, we utilize an \emph{absmean} quantization approach. This technique normalizes the weight matrix by dividing it by its mean absolute value, followed by mapping each entry to the closest integer within \{-1, 0, +1\}:

\begin{equation}
    \widetilde{W} = \mathrm{RoundClip}(\frac{W}{\gamma + \epsilon}, -1, 1),
\end{equation}

\begin{equation}
    \mathrm{RoundClip}(x, a, b) = \max(a, \min(b, \mathrm{round}(x))),
\end{equation}

\begin{equation}
    \gamma = \frac{1}{nm}\sum_{ij} |W_{ij}|.
\end{equation}

For activations, the quantization mechanism follows the scheme set in BitNet, but with a key difference: activations are not pre-scaled to $[0, Q_b]$ before applying nonlinearities. Instead, each activation is normalized to the range $[-Q_b, Q_b]$ on a per-token basis, thereby omitting zero-point quantization. This approach streamlines both the implementation and system optimization, while our results show it incurs almost no loss in accuracy.

\paragraph{LLaMA-alike Components.}

With LLaMA~\cite{llama, llama2} now serving as the primary template for open-source LLM development, our design of \text{BitNet b1.58} integrates LLaMA-inspired elements to foster compatibility within the open-source ecosystem. In particular, we incorporate RMSNorm~\cite{rmsnorm}, SwiGLU~\cite{swiglu}, and rotary embeddings~\cite{rope}, and eliminate all bias terms. This architectural choice facilitates seamless integration with widely adopted open-source libraries such as Huggingface, vLLM~\cite{vllm}, and llama.cpp\footnote{\url{https://github.com/ggerganov/llama.cpp}}, requiring only minimal adaptation.

\section{Results}

We conducted a comparison between \text{BitNet b1.58} and our implementation of FP16 \text{LLaMA LLM} across different parameter scales. Both sets of models were pre-trained from scratch on the RedPajama dataset~\cite{redpajama}, using 100 billion tokens, to guarantee a level playing field. We assessed their zero-shot capabilities on diverse language benchmarks, namely ARC-Easy~\cite{arc}, ARC-Challenge~\cite{arc}, Hellaswag~\cite{hellaswag}, Winogrande~\cite{winoGrande}, PIQA~\cite{piqa}, OpenbookQA~\cite{openbookqa}, and BoolQ~\cite{boolq}. Validation perplexity was also determined using WikiText2~\cite{wikitext2} and C4~\cite{c4} datasets.

For efficiency metrics, we measured and compared GPU memory usage and inference latency for both \text{LLaMA LLM} and \text{BitNet b1.58}. Evaluations were performed with the FasterTransformer\footnote{\url{https://github.com/NVIDIA/FasterTransformer}} platform, optimized for low-latency LLM inference on GPUs, with the inclusion of the 2-bit kernel from Ladder~\cite{ladder} tailored to \text{BitNet b1.58}. We provide the inference time required per output token, as this metric dominates overall serving cost.

Table~\ref{tab:ppl} presents a comparative overview of perplexity and resource requirements for \text{BitNet b1.58} and \text{LLaMA LLM}. The results indicate that, beginning at the 3B scale, \text{BitNet b1.58} reaches the perplexity performance of the FP16 \text{LLaMA LLM}, while demonstrating a 2.71-fold speedup and consuming 3.55 times less GPU memory. Specifically, the 3.9B configuration of \text{BitNet b1.58} is 2.4 times faster and reduces memory needs by a factor of 3.32, all while outperforming the \text{LLaMA LLM} at the 3B size in terms of language modeling accuracy.

Table~\ref{tab:zero-shot} details zero-shot accuracy across evaluation tasks. Utilizing the \emph{lm-evaluation-harness}\footnote{\url{https://github.com/EleutherAI/lm-evaluation-harness}} pipeline, our analyses demonstrate that the disparity in task performance between \text{BitNet b1.58} and \text{LLaMA LLM} lessens as model capacity grows. Notably, from the 3B size onwards, \text{BitNet b1.58} attains parity with its full-precision counterpart. This trend matches our observations on perplexity: \text{BitNet b1.58} at 3.9B not only surpasses the 3B \text{LLaMA LLM} in downstream metrics but does so with significantly reduced memory and latency. These findings highlight that \text{BitNet b1.58} constitutes a Pareto improvement over leading large language models.

\paragraph{Memory and Latency}

To investigate scalability, we expanded the model sizes to 7B, 13B, and 70B parameters and examined their associated costs. As depicted in Figure~\ref{fig:latency-memory-energy}, both latency and memory usage exhibit clear trends: larger models benefit from greater acceleration, with \text{BitNet b1.58} 70B demonstrating a 4.1-fold reduction in inference time relative to the \text{LLaMA LLM} baseline. This acceleration arises from the fact that the computational burden of \emph{nn.Linear} increases as the model size grows. Memory requirements similarly reflect these scaling properties, as the embeddings, which remain in full precision, represent a smaller proportion of total memory consumption for larger models. Both metrics were obtained using a 2-bit kernel, suggesting there remains significant potential for further cost reduction through additional optimizations.

\paragraph{Energy}

In addition, we quantified the energy required for arithmetic operations performed by both \text{BitNet b1.58} and \text{LLaMA LLM}. Our primary focus was on energy usage in matrix multiplications, given its dominant role in the computational cost of LLMs. Figure~\ref{fig:energy} presents the breakdown of total energy consumption: most of the energy in \text{BitNet b1.58} is attributed to INT8 addition operations, while \text{LLaMA LLM} relies on both FP16 addition and multiplication. Leveraging the energy estimation approach described in~\cite{energycost, pokebnn}, we found that \text{BitNet b1.58} reduces matrix multiplication energy consumption by a factor of 71.4 on 7nm hardware compared to the baseline. We further evaluated the overall energy consumed per inference for models processing 512 tokens, observing that as the models grow in size, \text{BitNet b1.58} increasingly outperforms the FP16 \text{LLaMA LLM} baseline in terms of energy efficiency. This pattern is attributed to the expanding influence of \emph{nn.Linear} operations at larger model scales, while the relative impact of other computational components diminishes.

\paragraph{Throughput}

To assess throughput, we evaluated both \text{BitNet b1.58} and \text{LLaMA LLM} models with 70B parameters, deploying them on a pair of 80GB A100 GPUs. Pipeline parallelism~\cite{gpipe} was utilized, enabling the execution of \text{LLaMA LLM} 70B on these GPUs. The batch size was incrementally increased until the GPU memory was saturated, maintaining a sequence length of 512 throughout the process. According to Table~\ref{tab:throughput}, \text{BitNet b1.58} 70B sustains up to 11 times larger batch sizes compared to \text{LLaMA LLM}, yielding a throughput improvement of 8.9-fold.

\textbf{\text{BitNet b1.58} introduces a novel scaling paradigm in relation to inference cost and model efficacy}. Drawing on observations from Figure~\ref{fig:latency-memory-energy} and Figure~\ref{fig:energy}, we provide the following equivalence between the capacities of 1.58-bit and 16-bit models:

\begin{itemize}
    \item A 13B BitNet b1.58 model demonstrates superior efficiency—spanning latency, memory, and energy metrics—over a 3B FP16 LLM.
    \item The 30B BitNet b1.58 model surpasses the 7B FP16 LLM regarding latency, memory demands, and energy use.
    \item A 70B BitNet b1.58 model outperforms a 13B FP16 LLM with respect to latency, memory footprint, and energy consumption.
\end{itemize}

\paragraph{Training with 2T Tokens}

Training token count remains a significant determinant for LLM performance. To investigate \text{BitNet b1.58}'s scalability relative to token quantity, we trained it on 2T tokens following the StableLM-3B~\cite{StableLM-3B-4E1T} data protocol—the current benchmark for open-source 3B models. We then assessed both models using a suite of benchmarks, including Winogrande~\cite{winoGrande}, PIQA~\cite{piqa}, SciQ~\cite{sciq}, LAMBADA~\cite{lambada}, and ARC-easy~\cite{arc}, reporting zero-shot accuracy in Table~\ref{tab:2t}. For benchmarks measured by both accuracy and normalized accuracy, their mean was computed. The results for StableLM 3b at the 2T-token scale are sourced directly from its technical documentation. Our evaluation indicates that \text{BitNet b1.58} surpasses its counterpart across all tested tasks, suggesting that 1.58-bit LLMs exhibit robust generalization performance.

\section{Discussion and Future Work}

\textbf{1-bit Mixture-of-Experts (MoE) LLMs}

Mixture-of-Experts (MoE) architectures offer an efficient solution for deploying LLMs by minimizing computational FLOPs, but their practical utility is still hampered by substantial memory requirements and the associated inter-chip communication burdens. The introduction of 1.58-bit LLMs has the potential to resolve these limitations. By dramatically lowering the memory footprint, fewer hardware devices are necessary for deploying MoE-based models. Furthermore, the volume of data transfer for activations across interconnected hardware is significantly decreased. Ideally, with sufficient compression, the entire model could reside on a single device, eliminating these communication bottlenecks altogether.

\textbf{Native Support of Long Sequence in LLMs}

Handling extended input sequences is becoming increasingly essential for LLM applications. A central obstacle in this context is the considerable memory usage demanded by KV cache storage during long sequence inference. \text{BitNet b1.58} constitutes an important advancement by shrinking activation precision from 16 bits down to 8 bits, which in turn permits the handling of sequences twice as long on the same hardware resources. Moreover, the activations may be further losslessly compressed, potentially reaching 4 bits or less with 1.58-bit LLMs—a direction we propose to investigate in subsequent studies.

\textbf{LLMs on Edge and Mobile}

The implementation of 1.58-bit LLMs can revolutionize the efficiency of large language models on edge and mobile hardware, both of which are frequently constrained in terms of memory and processing capabilities. Such hardware limitations have traditionally curtailed the feasible scale and performance of LLM deployments. However, by utilizing 1.58-bit LLMs’ lower energy consumption and reduced memory demands, deployment on these devices becomes more practical, paving the way for a host of new applications. This not only broadens the functional range of edge and mobile devices but also encourages novel use cases for LLMs in constrained settings. Additionally, since 1.58-bit LLMs are inherently more compatible with CPU architectures—the primary processing units in edge and mobile platforms—\text{BitNet b1.58} can execute efficiently, thereby maximizing device performance and functionality.

\textbf{New Hardware for 1-bit LLMs}

Emerging projects such as~Groq\footnote{\url{https://groq.com/}} have underscored the promise of developing custom hardware, like LPUs, tailored for LLM workloads. Building on this trend, we advocate for the development of specialized hardware and system-level designs expressly optimized for 1-bit LLM computation, as introduced by BitNet~\cite{bitnet}, to fully leverage this transformative approach.